% ===== 全反射・臨界角（原本 p117 の記述を加筆して丁寧にした） =====
% 元 ch24 は wrapfigure[6]{60mm} だったが，図の実高が46mm＝約9行分あり本文と重なっていた．
% zuright（minipage 2段）に置き換え，あわせてまとめ枠を新設した．
\hspace{1zw}屈折の法則$n_1\sin i=n_2\sin r$を屈折角について解くと$\sin r=\dfrac{n_1}{n_2}\sin i$となる．
そのため，{\bf 屈折率の大きな媒質から小さな媒質へ光が進む場合}（$n_1>n_2$）には$\sin r>\sin i$，すなわち{\bf 屈折角$r$の方が入射角$i$より大きくなる}．
このとき入射角$i$を少しずつ大きくしていくと屈折角$r$はそれより速く大きくなり，やがて$r=90^\circ$に達してしまう．
そこからさらに入射角を大きくすると{\bf 屈折光は存在できなくなり，入射光のすべてが境界面で反射される．}
この現象を{\bf 全反射}といい，全反射が始まる境目の入射角$\theta_0$を{\bf 臨界角}という．

\begin{zuright}{58mm}
\hspace{1zw}簡単のため，図のように絶対屈折率$n$の媒質中の点光源から出た光が真空中へと出ていく場合について考える．
入射角を$i$（媒質側），屈折角を$r$（真空側）とすると，真空の絶対屈折率は1であるから屈折の法則より
\[n\sin i=1\cdot\sin r\hspace{4mm}\text{すなわち}\hspace{4mm}\sin r=n\sin i\]
である．ここで$\sin r$は{\bf どんな角$r$に対しても1を超えられない}から，
\[\sin r=n\sin i>1\]
となるような入射角$i$に対しては，{\bf この式を満たす屈折角$r$がそもそも存在しない．}
\zusep
\includegraphics[keepaspectratio,width=58mm]{text_p117_fig1.png}
\end{zuright}

\hspace{1zw}屈折光が存在しないということは，境界面に達した光が行き先を失うということであり，このとき光は{\bf すべて反射される}．
すなわち，
\[\mbox{\boldmath $ \sin i>\dfrac{1}{n}$}\]
なる入射角$i$の光はすべて全反射することになる．
したがって臨界角$\theta_0$は，ちょうど屈折角が$r=90^\circ$となる境目の入射角であり，
\[n\sin\theta_0=1\cdot\sin 90^\circ\hspace{4mm}\Longleftrightarrow\hspace{4mm}\sin\theta_0=\dfrac{1}{n}\]
を満たす．

\begin{matome}{全反射と臨界角}
\hspace{3mm} \MARU{1} 全反射は{\bf 屈折率の大きい媒質から小さい媒質へ}進むときにしか起こらない．\\
\hspace{3mm} \MARU{2} $n_1>n_2$のとき，臨界角$\theta_0$は$\ \sin\theta_0=\dfrac{n_2}{n_1}\ $を満たし，$i>\theta_0$の光はすべて反射される．\\
\hspace{3mm} \MARU{3} {\bf 臨界角を求めるには，屈折の法則において屈折角を$90^\circ$と置けばよい．}\\
\hspace{3mm} 　　公式として覚える必要はない．「屈折角が存在しない条件」として毎回立式すればよい．
\end{matome}
